\section{A universal normalized ordinary ring for the cubic residual tuple}
\label{sec:cubic-residual-tuple-space}

The residual condition in the literal source \path{NRep.IsBar} is
\(\Psi^{w_i}\), with \(w=(1,0,-1)\) in rank three.
The cubic tuple does not satisfy that condition.  The construction
below replaces that condition explicitly and proves the required
universal ring, rather than applying the old \path{Space} to a
different residual representation.  No numerical \path{Frame.Ok},
\path{Big}, or deformation-dimension inequality is assumed here.

\paragraph{Data and the exact functor.}
Let \(\mathcal O\) be a complete discrete valuation ring with
finite residue field \(k\) of characteristic \(p\).
Let \(G\) be a profinite group with the following finiteness
condition, the literal source \path{PhiCond G p}:
\[
 \#\operatorname{Hom}_{\rm cont}(U,\mathbb F_p)<\infty
 \quad\text{for every open subgroup }U\subset G.
 \tag{G94.1}
\]
Let \(\star:G\to G\) be a continuous anti-involution, and
fix \(g_0\in G\).  For a finite set \(\mathcal V\), let
\(\iota_v:G_v\to G\) be continuous homomorphisms of profinite
groups.  Suppose given continuous characters
\(\theta_0,\theta_1,\theta_2:G\to k^\times\) satisfying
\[
 \theta_i(g^\star)=\theta_{2-i}(g).
 \tag{G94.2}
\]
For every \(v\), require an element \(s_v\in G_v\) such that
\[
 \theta_0(\iota_v s_v),\quad
 \theta_1(\iota_v s_v),\quad
 \theta_2(\iota_v s_v)
 \quad\text{are pairwise distinct}.
 \tag{G94.3}
\]
The finite images of these characters make their diagonal
representation \(\bar r\) locally constant.

The test rings \(C\) are complete Noetherian local
\(\mathcal O\)-algebras with residue field identified with
\(k\) through \(\mathcal O\), exactly the source
\path{IsCO O C} category.  A normalized ordinary lift over
\(C\) consists of a continuous representation
\(r:G\to\operatorname{GL}_3(C)\), reducing to
\(\bar r=\operatorname{diag}(\theta_0,\theta_1,\theta_2)\),
and continuous local characters
\(\kappa_{v,i}:G_v\to C^\times\), reducing to
\(\theta_i\iota_v\), with
\[
 \begin{split}
 r(g^\star)&=Jr(g)^tJ,\\
 \bigl(r(g_0)-r(g_0^\star)\bigr)_{ij}&=0\quad(i\ne j),\\
 \prod_{i=0}^{2}
   \bigl(r(\iota_v h_i)-\kappa_{v,i}(h_i)I\bigr)&=0
       \quad(h_0,h_1,h_2\in G_v).
 \end{split}
 \tag{G94.4}
\]
Here \(J\) is the antidiagonal matrix with entries one, and
the last product is in the displayed order.  Lifts are in the
specified basis: there is no quotient by strict conjugation.
In particular no scalar-endomorphism or residual-irreducibility
hypothesis is being imposed.

\begin{theorem}[AC-G94-space: the changed-residual universal ring]
\label{AC-G94-space}
Under (G94.1)--(G94.3), the functor (G94.4) is represented by
a nonzero complete Noetherian local \(\mathcal O\)-algebra
\(R_{\theta}\) with residue field \(k\).
More precisely, there are a universal pair
\((r_{\theta},\kappa_{\theta})\) and, for every test ring
\(C\), a natural bijection
\[
 \operatorname{Hom}_{\mathcal O\text{-alg}}(R_{\theta},C)
 \simeq
 \{\text{normalized ordinary lifts in (G94.4) over }C\}.
 \tag{G94.5}
\]
Every homomorphism on the left is automatically local and
continuous.  The local characters on the right are uniquely
determined by the representation.

This is the literal normalized ordinary universal-property
construction with \path{IsBar} replaced by the specified
tuple.  It supplies neither a dimension formula nor the
\(A/B\), \path{FPt}, or \path{At} comparison for this new ring.
\end{theorem}

\begin{proof}
\emph{1. Start with the genuinely generic framed lift ring.}
The source theorem
\path{CommutativeAlgebra/UnivLift.lean}, \path{univ_lift},
applies to an arbitrary locally constant residual matrix
representation of a profinite group satisfying (G94.1).
Its conclusion gives a complete Noetherian local ring
\(R^{\square}\), residue field \(k\), and a universal
continuous lift \(r^{\square}\) of \(\bar r\).
This is the one general representability input retained here.
The proof of the containing source file was inspected in full;
its finite tangent-separation and successive Artin-level
construction do not require \(\bar r\) to be irreducible.
There is no use at this step of a residual character and its
powers.

We record also why the maps in (G94.5) need no separate
continuity clause.  For an \(\mathcal O\)-algebra map
\(f:A\to C\) between these local algebras and \(x\in\mathfrak m_A\),
suppose its residue \(\overline{f(x)}\) were nonzero.
Choose \(a\in\mathcal O\) whose residue is its inverse.
Then \(1-ax\) is a unit in \(A\), whereas its image has zero
residue in \(C\), a contradiction.  Thus
\(f(\mathfrak m_A)\subset\mathfrak m_C\), the residue map is
the prescribed identity, and
\(f(\mathfrak m_A^n)\subset\mathfrak m_C^n\) proves continuity.

\emph{2. Diagonalize one local separating element integrally.}
For each \(v\), put
\(M_v=r^{\square}(\iota_v s_v)\).
Its residual matrix is diagonal with distinct entries.
We give the Hensel step explicitly.
The characteristic polynomial of \(M_v\) reduces to the
product of three distinct linear factors.  Simple-root
Hensel lifting gives uniquely labeled roots
\(d_{v,i}\in R^{\square}\), lifting
\(\theta_i(\iota_v s_v)\).  Their differences are units,
and the characteristic polynomial is their product.
For instance, divide successively by \(X-d_{v,i}\);
the unit differences imply that the other roots remain
roots of each successive quotient, proving this factorization.

The Lagrange projectors
\[
 E_{v,i}=\prod_{j\ne i}
       \frac{M_v-d_{v,j}I}{d_{v,i}-d_{v,j}}
 \tag{G94.6}
\]
are orthogonal idempotents, sum to one, and satisfy
\(M_vE_{v,i}=d_{v,i}E_{v,i}\).  These statements follow by
polynomial interpolation modulo the characteristic polynomial
and Cayley--Hamilton.  Their reductions are the coordinate
projectors.  Let the \(i\)-th column of \(P_v\) be
\(E_{v,i}e_i\), divided by its \(i\)-th coordinate, which
is a unit.  Then
\[
 P_v\equiv I\pmod{\mathfrak m_{R^{\square}}},\qquad
 P_v^{-1}M_vP_v=\operatorname{diag}(d_{v,0},d_{v,1},d_{v,2}).
 \tag{G94.7}
\]
Invertibility follows already from the reduction of its
determinant.  This construction commutes with maps of the
test rings: Hensel roots have the fixed residual labels,
and every denominator in (G94.6) remains a unit.
Only this single matrix is diagonalized; arbitrary local
deformations have not been assumed to be split.

Write
\[
 T_v(h)=P_v^{-1}r^{\square}(\iota_v h)P_v.
 \tag{G94.8}
\]
Its entries are continuous, and its reduction is the
prescribed diagonal tuple for every \(h\).

\emph{3. Impose an explicit closed ideal.}
Let \(I_{\theta}\) be the ideal of \(R^{\square}\)
generated by all entries of
\[
 \begin{gathered}
 r^{\square}(g^\star)-Jr^{\square}(g)^tJ,\\
 \bigl(r^{\square}(g_0)-r^{\square}(g_0^\star)\bigr)_{ij}
       \quad(i\ne j),\\
 T_v(h)_{ab}\quad(b<a).
 \end{gathered}
 \tag{G94.9}
\]
All these elements reduce to zero.  For the first group
this is (G94.2), for the second it is diagonality, and
for the third it is the reduction in (G94.8).
Consequently \(I_{\theta}\subset\mathfrak m_{R^{\square}}\).
Set
\[
                 R_{\theta}=R^{\square}/I_{\theta}.
 \tag{G94.10}
\]
It is nonzero and has residue field \(k\).
Although (G94.9) is indexed by group elements, it generates
an ordinary ideal in a Noetherian ring.  All such ideals
are closed in its maximal-ideal topology; equivalently the
source \path{co_quot} applies.  The quotient is therefore
complete and Noetherian.  No flatness over \(\mathcal O\)
or absence of vertical torsion is asserted.

Over this quotient every \(T_v(h)\) is upper triangular.
Its diagonal entries
\[
                  \kappa_{\theta,v,i}(h)=T_v(h)_{ii}
 \tag{G94.11}
\]
are continuous multiplicative characters.  They are units:
the diagonal of \(T_v(h)T_v(h^{-1})=I\) gives their
inverses.  Their residual values are the required ones.
For upper triangular matrices the ordered product in
(G94.4) is zero.  Indeed its first factor has zero first
diagonal entry and hence kills the first column; after
multiplication by the second factor the first two columns
vanish; after the third all three vanish.  Conjugating by
\(P_v\) proves the ordered-product identity before the
change of basis.  The other two conditions in (G94.4)
were imposed directly.  This constructs the universal pair.

\emph{4. The ordered-product equations force the imposed ideal.}
Let \((r,\kappa)\) be any pair (G94.4) over a test ring
\(C\).  Generic framed universality gives a unique
\(f:R^{\square}\to C\) carrying \(r^{\square}\) to \(r\).
Put \(P=f(P_v)\),
\(D=\operatorname{diag}(f(d_{v,0}),f(d_{v,1}),f(d_{v,2}))\),
and \(T(h)=P^{-1}r(\iota_v h)P\).
Then \(T(s_v)=D\).
The ordered-product identities are unchanged by this
constant conjugation.

Fix \(a\in\{0,1,2\}\), put \(h_a=h\), and put
\(h_j=s_v\) for \(j\ne a\).
The \((a,b)\)-entry of the resulting identity, for
\(b\leq a\), is
\[
 \left(\prod_{j<a}(D_{aa}-\kappa_{v,j}(s_v))\right)
 \bigl(T(h)-\kappa_{v,a}(h)I\bigr)_{ab}
 \left(\prod_{j>a}(D_{bb}-\kappa_{v,j}(s_v))\right)=0.
 \tag{G94.12}
\]
Every factor outside the middle entry is a unit.
On the left \(j<a\) excludes the residual label \(a\);
on the right \(j>a\geq b\) excludes the residual label
\(b\).  Condition (G94.3) therefore proves this unit claim
without any reducedness assumption on \(C\).
It follows that
\[
 T(h)_{ab}=0\ (b<a),\qquad
 T(h)_{aa}=\kappa_{v,a}(h).
 \tag{G94.13}
\]
Thus \(f\) kills the third family in (G94.9), and the
first two families vanish by the supplied polarization
and normalization.  It factors through (G94.10).
The diagonal equalities (G94.13) identify its local
characters with (G94.11).  They also show that the
characters were uniquely determined by \(r\).
Finally any two maps from \(R_{\theta}\) inducing this
pair have the same composite from \(R^{\square}\), by
framed universality, and hence are equal.  This proves
(G94.5) and its naturality.
\end{proof}

\paragraph{Actual arithmetic finiteness and the cubic specialization.}
There is no new arithmetic \(\Phi\)-condition hidden in this
construction.  If \(\Gamma\) is the actual G90 quotient for
a number field and a finite residual splitting field, choose
\(S\) containing the coefficient-prime places and every place
where that splitting field ramifies.  Its defining away-kernel
contains all inertia outside \(S\).  Thus \(\Gamma\), and
its relevant finite-index subgroups, are quotients of the
corresponding unramified-outside-\(S\) Galois groups.
Those groups satisfy (G94.1).  Here is the precise finiteness
argument.  For any associated finite number field \(B\),
restriction to \(B(\mu_p)\) is injective on
\(H^1(-,\mathbb F_p)\), since its degree is prime to \(p\).
After choosing a primitive \(p\)-th root of unity, Kummer
theory identifies this group with a subgroup of the group
in the exact sequence
\[
 0\longrightarrow
 \mathcal O_{B(\mu_p),S}^{\times}/
       (\mathcal O_{B(\mu_p),S}^{\times})^p
 \longrightarrow H^1(G_{B(\mu_p),S},\mu_p)
 \longrightarrow\operatorname{Cl}(\mathcal O_{B(\mu_p),S})[p]
 \longrightarrow0.
 \tag{G94.14}
\]
Both outer groups are finite.  Inflation from a quotient
is injective, so this proves (G94.1) for every open subgroup
of the actual quotient as well.  It uses only the standard
Kummer, finite class-group and finitely generated \(S\)-unit
inputs already retained in the arithmetic audits.

For the actual cubic data on the higher host of AC-G93-fill,
take
\[
 (\theta_0,\theta_1,\theta_2)
       =(\bar\chi^j\gamma,1,\bar\chi^j\gamma^{-1}),
 \qquad g^\star=cg^{-1}c^{-1}.
 \tag{G94.15}
\]
The rational character \(\gamma\) is fixed by \(c\), and
\(\bar\chi^c=\bar\chi^{-1}\).  Hence
\(\bar\chi(g^\star)=\bar\chi(g)\) and
\(\gamma(g^\star)=\gamma(g)^{-1}\), proving (G94.2).
At each actual coefficient-prime local map, the retained
separation provides an inertia element with
\(\gamma^2\ne1\).  Cubing kills the cubic character while
preserving that inequality, because the order of this
inertia value divides \(268\), which is prime to three.
The three local values are then \((t,1,t^{-1})\), with
\(t^2\ne1\).  All these residual images have order prime
to \(269\), so passage to the selected-place and cyclotomic
\(269\)-power extensions preserves the same separating
images.  This proves (G94.3) in the actual range.
The desired \(g_0\) is the normalization witness already
constructed in AC-G90; no additional separation at
\(g_0\) is needed merely to construct this ring.

\paragraph{Literal source scope and the remaining interface.}
The complete containing source files read for this result are
\path{FrobeniusSmoothVectors/SpaceEx.lean},
\path{LeAdicClosure.lean}, \path{OrdOfFlag.lean},
\path{AlgebraTools/MonomialRepresentations/HenselDiag.lean},
\path{CommutativeAlgebra/UnivLift.lean},
\path{CoQuot.lean}, and
\path{PatchingModules/DeterminantPadicFree/PhiSub.lean},
at the pinned source commit recorded in the accompanying manifest.
The generic representability theorem is retained with its full
literal hypotheses.  The quotient, Hensel diagonalization,
ordered-product equivalence and arithmetic specialization are
proved above.  No Lean execution is claimed.

The exact reusable part of \path{NRep} is its representation,
polarization, normalization and ordered-product equations.
Its map/restrict constructions are polynomial and unchanged.
The exact changed part is \path{IsBar}, and consequently
the residual clauses in \path{IsLift}, \path{Integral},
\path{Cone}, and \path{Space}.  This theorem supplies the
new \path{Space} universal property; it does not identify
its graded subalgebras or characteristic-fiber points with
the original single-character objects.  In particular the
\path{HostFrameEx} abelian/single-kernel input, the source
numerical inequalities, and prescribed Hecke-curve support
remain separate requirements.

\paragraph{Geometric dictionary.}
The formal spectrum of \(R_{\theta}\) parametrizes framed
deformations of the fixed residual lisse local system, with
the stated conjugate-duality normalization and ordered local
flags.  The separating element makes the local flag equations
visible in one integral basis.  This uses the previously audited
finite-cover and deformation-ring dictionary; it asserts no
new archival SGA reading and no automorphy of every point.
